\section{Desenvolvimento e Implementação}

\subsection{Arquitetura Modular em ROS 2 Jazzy e Gazebo Harmonic}
% TODO: Descrever os pacotes, plugins SDF do modelo robo_diferencial (diff-drive, joint-state-publisher, gpu_lidar, odometry-publisher), bridge ros_gz_bridge (bridge_params.yaml), QoS (qos_profile_sensor_data em /scan) e sincronizacao temporal (/clock).
% Referenciar a Figura~\ref{fig:arquitetura_ros2} e as referencias \cite{ros2jazzy,gazeboharmonic}.

\begin{figure}[htbp]
\centering
\resizebox{\columnwidth}{!}{\begin{tikzpicture}[
    font=\scriptsize,
    >=Stealth,
    sim/.style={
        rectangle, rounded corners=2.5pt, draw=black!80, fill=gray!12,
        thick, minimum width=3.3cm, minimum height=0.95cm, align=center, inner sep=4pt
    },
    bridge/.style={
        rectangle, rounded corners=2.5pt, draw=teal!75!black, fill=teal!10,
        thick, minimum width=3.3cm, minimum height=0.95cm, align=center, inner sep=4pt
    },
    rosnode/.style={
        rectangle, rounded corners=2.5pt, draw=blue!70!black, fill=blue!8,
        thick, minimum width=3.7cm, minimum height=0.95cm, align=center, inner sep=4pt
    },
    eval/.style={
        rectangle, rounded corners=2.5pt, draw=orange!80!black, fill=orange!10,
        thick, minimum width=3.7cm, minimum height=0.95cm, align=center, inner sep=4pt
    },
    topic/.style={
        font=\ttfamily\scriptsize, inner sep=1.5pt
    }
]

% Coluna esquerda: Simulador e Ponte
\node[sim]    (gz)     at (0, 0.3)    {\textbf{Gazebo Harmonic}\\\texttt{robo\_diferencial} (.sdf)};
\node[bridge] (br)     at (0, -1.7)   {\textbf{\texttt{ros\_gz\_bridge}}\\\texttt{/clock}, sensores e \texttt{/cmd\_vel}};

% Coluna direita: Avaliacao, Estimacao e Controle
\node[eval]    (logger) at (7.2, 0.3)  {\textbf{Nós Avaliadores}\\\texttt{metrics\_logger} / \texttt{nav\_metrics}};
\node[rosnode] (odom)   at (7.2, -1.7) {\textbf{\texttt{odom\_estimator}}\\Integração RK2 ($50\text{ Hz}$)};
\node[rosnode] (nav)    at (7.2, -3.7) {\textbf{Controlador Reativo}\\\texttt{apf\_navigator} / \texttt{bug2\_navigator}};

% 1) Gazebo <-> Bridge
\draw[<->, thick] (gz.south) -- node[right=2pt, font=\scriptsize] {Gz Transport} (br.north);

% 2) Ground Truth isolado (Bridge -> Avaliadores)
\draw[->, dashed, thick, orange!85!black]
    ([yshift=0.22cm]br.east) -- (2.05, -1.48) -- (2.05, 0.3)
    -- node[midway, above=2pt, topic, text=orange!75!black] {/ground\_truth/odom} (logger.west);

% 3) Encoders (Bridge -> Odom)
\draw[->, thick]
    ([yshift=-0.12cm]br.east) -- ([yshift=-0.12cm]odom.west)
    node[pos=0.58, below=2pt, topic] {/joint\_states};

% 4) LiDAR (Bridge -> Controlador)
\draw[->, thick]
    ([xshift=0.55cm]br.south) |- ([yshift=0.12cm]nav.west)
    node[pos=0.73, above=2pt, topic] {/scan (best\_effort)};

% 5) Odometria estimada (Odom -> Controlador e Odom -> Avaliadores)
\draw[->, thick] (odom.south) -- node[right=2pt, topic] {/odom\_est} (nav.north);
\draw[->, thick] (odom.north) -- node[right=2pt, topic] {/odom\_est} (logger.south);

% 6) Comando de velocidade (Controlador -> Bridge)
\draw[->, thick]
    (nav.south) -- ++(0, -0.55) -| ([xshift=-0.55cm]br.south)
    node[pos=0.25, below=2pt, topic] {/cmd\_vel};

% 7) Status e eventos (Controlador -> Avaliadores pela direita)
\draw[->, thick]
    (nav.east) -- ++(0.45, 0) |- (logger.east)
    node[pos=0.25, right=2pt, topic, align=left] {/nav\_status\\/nav\_events};

\end{tikzpicture}
}
\caption{Arquitetura modular em ROS~2 Jazzy e Gazebo Harmonic com isolamento estrito do tópico \texttt{/ground\_truth/odom}.}
\label{fig:arquitetura_ros2}
\end{figure}

\subsection{Estimador de Odometria (\texttt{odom\_estimator\_node})}
% TODO: Detalhar a implementacao do no odom_estimator_node.py e do modulo kinematic_model.py (leitura de /joint_states, injecao opcional de ruido gaussiano com seed fixa, passo RK2 e publicacao em /odom_est).

\subsection{Controladores Reativos: APF e Bug2}

Ambos os controladores operam em malha fechada a $20\text{ Hz}$ assinando exclusivamente \texttt{/odom\_est} e \texttt{/scan}, sem acesso ao \textit{ground truth}, e encerram a navegação com sucesso quando a distância estimada até a meta é inferior à tolerância $d_{\text{tol}}$ ($0{,}20\text{ m}$).

No nó \texttt{apf\_navigator\_node}, a força resultante é obtida somando a atração para a meta às contribuições repulsivas dos feixes do LiDAR que detectam obstáculos dentro do raio $\rho_0$. Como o robô diferencial não se desloca lateralmente, a direção e a magnitude dessa força resultante são convertidas em comandos de velocidade angular e linear por meio dos ganhos cinemáticos $k_\omega$ e $k_v$ (Tabela~\ref{tab:parametros}), reduzindo a velocidade linear quando o erro de orientação é elevado para que o robô gire antes de avançar. Quando o deslocamento do robô permanece desprezível por um intervalo contínuo, o controlador identifica automaticamente a estagnação por mínimo local ou GNRON.

No nó \texttt{bug2\_navigator\_node}, o algoritmo clássico~\cite{lumelsky1987} foi implementado como uma máquina de estados finitos alternando entre busca direta (\texttt{GOAL\_SEEKING}) e contorno de obstáculo (\texttt{WALL\_FOLLOWING}), com ajustes práticos necessários para lidar com as dimensões físicas do robô, a discretização temporal e o ruído sensorial~\cite{siegwart2004}. No modo \texttt{GOAL\_SEEKING}, o robô segue a linha-$m$ rumo à meta até detectar um obstáculo frontal dentro do limiar $d_{\text{hit}}$, registrando o ponto de impacto $H^j$ e iniciando o contorno. No modo \texttt{WALL\_FOLLOWING}, o controlador mantém o obstáculo a uma distância lateral segura $d_{\text{wall}}$ até satisfazer simultaneamente as condições práticas de saída para o ponto $L^j$: ter percorrido um trecho mínimo de contorno $d_{\text{debounce}}$ (evitando trocas falsas de estado logo após $H^j$), reencontrar a linha-$m$ dentro da faixa de tolerância $\delta_m$, estar mais próximo da meta do que $H^j$ por uma margem mínima $\epsilon$ e possuir caminho frontal livre na direção do alvo.

\subsection{Isolamento do \textit{Ground Truth} e Registro de Métricas}
% TODO: Explicar a garantia arquitetural de que apenas metrics_logger_node.py e nav_metrics_logger_node.py assinam /ground_truth/odom, alem da orquestracao automatizada das baterias (scripts run_etapa1.py e run_fase6_experimentos.py).
